\section{Repeat retained bits when expanding}\label{rgb-expansion}
Read the low and high input bytes unsigned, then combine them as low plus 256 times
high. Extract the five-bit red field, six-bit green field and five-bit blue field.
The expansion rule is
\[
 R'=8r+\left\lfloor r/4\right\rfloor,\qquad
 G'=4g+\left\lfloor g/16\right\rfloor,\qquad
 B'=8b+\left\lfloor b/4\right\rfloor.
\]
For a five-bit component, shift its retained bits left by three and copy its highest
three bits into the empty positions. For a six-bit component, shift left by two and
copy its highest two bits. The shifted ranges do not overlap, so OR implements
the same combination. Maximum fields expand to 255 rather than stopping at 248 or
252, as zero-filling would.

The worked packed value \texttt{7968} expands to $(123,44,66)$. Its original green
and blue low bits cannot be recovered from the packed bytes. Expanding the same
stored value always gives this result, regardless of which original input selected
that bucket. Packing the expanded values again recovers the same RGB565 value.
The reverse round trip generally changes an RGB888 input.

Bit repetition is the exact rule for this exercise. It is not defined as rounding
an ideal real-valued rescaling. Implementing a visually similar interpolation can
therefore disagree at particular field values even when black and white agree.
